\chapter{Pair Production}

The previous chapter perfected a derivative and, in perfecting it, told the story
slightly backwards. It celebrated the linear part of the variation --- the first
variation, exact and unique and with remainder identically zero --- as though the
zero remainder were the achievement. But a remainder that is identically zero is a
remainder that has been thrown away, and the thing thrown away is the prize. The
step that builds the full derivative \emph{adds} one datum, the residue; the step
that builds only the linear part \emph{projects that datum out}. The residue is the
real content of the computation; the linear part is merely the computation. This
chapter recovers the discarded datum, names it, and discovers that it is a signed
unit of charge.

The chapter begins with the correction itself. The earlier tower --- the
refinements, the squeeze, the produced certificate --- all worked in the weak form,
the form in which the residue is projected out before anyone looks. That is why the
residue stayed anonymous for so long. It was present at every stage, sitting in the
slot the weak form discards, and no instrument the book had built could see it. The
first section states the correction plainly: the order-two residue, not the
order-one linear part, is what the derivative computation produces, and the price of
working in the convenient weak form was exactly the loss of that residue.

The second section explains why the residue stayed hidden through every earlier
measurement, and the explanation is structural rather than accidental. A variation
that moves a single coordinate reads the mixed coupling as zero --- not approximately,
not usually, but always, for every action and every path. The cross term that pairs
the two sides of the middle cancels identically when only one side moves. So every
one-coordinate probe in the book, including the kicked path that read a clean value
of two, read the coupling as nothing. The instruments were blind to the residue by
construction, and the blindness was complete. A quantity that lives strictly between
two coordinates cannot be seen by any instrument that moves one coordinate at a time.

The third section produces the residue the only way it can be produced: in a pair. It
kicks both middle nodes of the flat vacuum at once, under the discriminating action,
and subtracts the story from the measurement. The two one-coordinate computations,
added together, predict that the action changes by four. The action actually changes
by three. The difference --- what is left after the linear story almost explains the
measurement --- is a signed unit. Read as measurement minus story it is negative one,
and the chapter names it the electron. Read in the opposite orientation, story minus
measurement, it is positive one, and that is the positron. The pair kick produces a
pair, signed, of unit magnitude, and produced only in coupling. The early chapters
left a pressure around ``what is left after the process almost explains itself,''
around ``the difference between the story and the measurement.'' In this section,
that pressure becomes an exact identity.

The fourth section ties the residue to the strain the book named at the outset. The
linear story fails to account for the measurement precisely when the residue is
nonzero; the failure and the residue are not two facts but one, an equivalence. The
strain the early chapters spoke of --- the informational strain that a variation
leaves behind --- is exactly the residue, and at the electron it is real because the
electron's charge is nonzero, and only because. The section seats this strain back
into the apparatus's structure: the residue question the foundational chapters left
open now has a precise slot, and the apparatus's own machinery manufactures exactly
the electron's term from the paired variation. The earlier pressure was there; this
chapter identifies the slot and fills it.

The fifth section certifies the configuration and reads the answer against a problem
a century old. A cubic spline's defining property is that it carries data through the
second derivative and nothing above --- it is twice continuously differentiable, what
the chapter calls \(C^2\). A pair of certificates, one per middle node, certifies the
configuration through order two: the sub-paths admit at order zero, the
Euler-Lagrange residuals vanish at order one, and what remains of every coupled
variation at order two is exactly the second variation, with the exactness leaving no
order three to certify. The order-two data of the certified spline is the residue
field, and the electron is its order-two quantum. Read against Hilbert's sixth
problem --- the call, in 1900, for a mathematical treatment of the axioms of physics
--- the section offers a finite variational mechanics with equations of motion at
order one and interaction at order two, exact at order two, derived as a theorem from
the apparatus's measurement primitives and machine-checked. It is one electron's
worth of physics, and the chapter claims exactly that much.

By the end of the chapter the single invariant of Chapter 8 has a physical face. The
order-two quantity that survived every first-order choice, that the apparatus could
read but not yet interpret, turns out to be a signed unit of charge, produced in
pairs, invisible to one-coordinate measurement, and certified through second order.
The book set out to measure a single invariant. It has measured the electron.

\section{The Correction to the Derivative Story}

The previous chapter's triumph needs one correction, and the correction reverses an
emphasis rather than a fact. The chapter built the first variation, proved it exact
and unique, and celebrated that its remainder was identically zero. Every theorem
there is true. But the celebration pointed at the wrong term. A derivative
computation, done honestly, produces two things: a linear part and a residue. The
chapter kept the linear part and discarded the residue, and then praised the
discarding as a clean result.

State the two steps the way the apparatus actually performs them. Building the full
derivative from a coarser variation \emph{adds} a datum --- the residue, the order-two
remainder that records what the linear part fails to capture. Building only the linear
part from the full derivative \emph{removes} that datum --- projects it out, keeps the
computation, loses what the computation was for. The earlier chapters lived on the
second step. They worked in the weak form, where the residue is projected away before
the measurement is read, and in that form the residue is structurally absent, not
merely small.

So the order-two residue, not the order-one linear part, is the derivative's real
content. Write the decomposition once more, but read it now with the emphasis
corrected:
\[
  D(v) = \underbrace{r_{\mathrm{L}}(v_{\mathrm{L}}) + r_{\mathrm{R}}(v_{\mathrm{R}})}_{\text{the story (linear)}}
         + \underbrace{\delta^2(v)}_{\text{the residue (order two)}}.
\]
The earlier chapters drove the story to zero and read the result as a solved
variational principle. That reading is correct and incomplete. When the story
vanishes, what is left is not nothing; it is the residue, the order-two datum the weak
form had been discarding all along. The solved principle does not eliminate the
residue. It isolates it.

This corrects the book's own language without retracting a single theorem. The
order-one residual of the previous chapter keeps its name and its uniqueness; it is a
genuine first-order quantity, exactly determined. The residue of this chapter is a
different slot, one floor up the ladder: the order-two datum that the first-order
computation cannot reach. The two are not competitors. The first-order residual is
the story; the order-two residue is what the story leaves out. The previous chapter
perfected the story. This chapter recovers what the story left out.

The cost of the weak form has a name the apparatus has used before: it is the price of
working with a projected, convenient computation rather than the full one. That price
bought every clean result of the tower arc --- the squeeze, the eventual exactness,
the produced certificate --- and it bought them by setting the residue aside. There is
nothing wrong with paying that price when the goal is to solve the first-order
equations. The error would be to forget that it was paid, and to mistake the projected
computation for the whole. The residue was never destroyed. It was set down in a slot
the weak form does not read, waiting for an instrument that could pick it up.

The section's metaphysical anchor is that what a computation discards can be its
point. The apparatus tanges the variation into a linear story and an order-two
residue, and the weak form funges the story forward while letting the residue fall
silent. The correction is not to redo the computation but to look at the slot the
computation emptied. The derivative's residue --- the order-two remainder, the second
variation, the single invariant under a new name --- has been present at every stage of
the book. The remaining chapters are about picking it up.

\section{Why One-Coordinate Instruments Are Blind}

The residue stayed anonymous for fifty chapters' worth of measurement, and the
section's task is to show that this was not an oversight. It was forced by the shape of
every instrument the book had built. Each of those instruments moved one coordinate at
a time, and a one-coordinate instrument reads the residue as exactly zero --- always,
for every action and every path. The blindness was structural and complete.

The residue is the mixed coupling, the cross term that pairs the left and right sides
of the middle. Look at what happens to that cross term when only one side moves. Move
the left coordinate and hold the right fixed at its vacuum value. The mixed coupling is
built from four cost readings that, with the right coordinate unmoved, cancel in two
pairs:
\[
  \delta^2(\text{left only}) = 0
  \qquad\text{for every action and every path.}
\]
The same cancellation happens when only the right coordinate moves:
\[
  \delta^2(\text{right only}) = 0
  \qquad\text{for every action and every path.}
\]
These are identities, not approximations. A one-coordinate variation does not read a
small coupling that could be refined away; it reads the coupling as precisely nothing,
because the cross term requires both coordinates to move and one of them is held still.

This explains the entire history of the book's measurements. Every probe the apparatus
ran moved one node at a time. The kicked path of the physical-record chapter, which
read a clean value of two, moved a single middle node; that two was entirely the
linear story, and the coupling underneath it read zero. The tower's refinements, the
squeeze, the produced certificate --- all of them exercised the apparatus one
coordinate at a time, and so all of them read the residue as zero and reported only the
story. The residue was present at every one of those stages, in the slot the weak form
discards, and invisible to every instrument that moved one thing at a time.

The blindness is worth stating as a principle, because it is the chapter's hinge. A
quantity that lives strictly in the coupling between two coordinates cannot be detected
by any instrument that varies one coordinate. No refinement helps, no matter how fine,
because refinement of a one-coordinate probe is still a one-coordinate probe. No change
of action helps, because the cancellation holds for every action. The only way to make
the residue nonzero is to move both coordinates at once, and an instrument built to
move one at a time can never do that. The residue is not hard to see; it is impossible
to see, with a single-coordinate instrument, by the algebra of the cross term.

There is a lesson here about measurement in general, and the apparatus states it. An
instrument's reach is bounded by the shape of the variations it can perform. An
instrument that probes along one axis is not merely less sensitive to a coupling
between axes; it is exactly, identically insensitive to it. Sensitivity is not a
matter of degree here but of structure: the coupling lies in a direction the
one-coordinate instrument cannot point. To measure it at all requires a different kind
of instrument --- one that moves a pair.

The section's metaphysical anchor is that an instrument can be blind by construction,
and that such blindness is not a flaw to be refined away but a fact to be respected.
The apparatus tanges a single coordinate into a reading and funges that reading
forward, and the cross term it cannot reach stays exactly zero in the record. The
residue's long anonymity was not the apparatus failing to look hard enough. It was the
apparatus looking, faithfully and finely, in a direction where the residue casts no
shadow at all. To see the residue, the next section stops moving one coordinate and
moves two.

\section{The Pair Kick and Measurement Minus Story}

To produce the residue, the apparatus moves a pair. It takes the flat vacuum --- the
path whose middle nodes sit at the trivial configuration --- and kicks both middle
nodes at once, under the discriminating action, sending the left node to one
representative and the right node to a different one. This is the pair kick, and it is
the first variation in the book that moves two coordinates simultaneously. What it
produces is the residue the one-coordinate instruments could never read.

Read the two quantities the pair kick presents. The story is the linear account: the
two one-coordinate computations, the left first-order residual and the right
first-order residual, added together. Under this pair kick the story totals four. The
measurement is the full coupled difference, the actual change in the action under the
joint variation. Under this pair kick the measurement totals three. The story and the
measurement disagree, and the disagreement is the whole point:
\[
  \text{story} = r_{\mathrm{L}} + r_{\mathrm{R}} = 4,
  \qquad
  \text{measurement} = D = 3.
\]
The residue is what is left when the story is subtracted from the measurement:
\[
  \delta^2(\text{pair kick}) = \text{measurement} - \text{story} = 3 - 4 = -1.
\]
A signed unit. The apparatus names it the electron. It has unit magnitude, negative
sign, and it is produced only in coupling --- a one-coordinate kick would have given
zero, by the previous section, and only the pair makes it appear.

This residue is not a new quantity. It is the second variation of the action,
evaluated at the pair kick:
\[
  \text{electron} = \delta^2(\text{pair kick}) = -1.
\]
The single invariant of Chapter 8 --- the order-two remainder that survived every
first-order choice, forced and unique --- evaluated on a paired excitation of the
vacuum, is a signed unit of charge. The book promised a single invariant and spent
four chapters making it readable; here the invariant, read on the simplest nontrivial
paired excitation, is the electron.

The sign is not a convention the apparatus may discard, and the chapter pays the debt
its title implies. The residue is signed, and the apparatus has chosen to read it as
measurement minus story, which makes the produced charge negative. The opposite
orientation --- story minus measurement --- reads the same magnitude with the opposite
sign:
\[
  \text{measurement} - \text{story} = -1 \ \ (\text{electron}),
  \qquad
  \text{story} - \text{measurement} = +1 \ \ (\text{positron}).
\]
Pair production produces a pair. The negative unit and the positive unit are the same
residue read in two orientations, and the apparatus cannot produce one without the
other being available as its mirror. The title is literal: the pair kick produces the
electron and the positron together, as the two orientations of one signed residue.

The magnitude depends on the pair being genuinely distinct. Kick both middle nodes to
the \emph{same} representative and the residue reads negative two, not negative one ---
a different excitation, not the unit charge. The unit magnitude belongs specifically
to the pair that moves the two coordinates to distinct values, the genuinely
two-channel excitation. The apparatus does not get a unit from any pair; it gets a
unit from the pair that excites both channels differently.

This production answers, exactly, an early residue question the book kept open. The
apparatus asked what remains when its story almost explains its measurement, before it
had machinery precise enough to carry the answer. That open question is now a theorem:
\[
  \delta^2(v) = D(v) - \big(r_{\mathrm{L}}(v_{\mathrm{L}}) + r_{\mathrm{R}}(v_{\mathrm{R}})\big).
\]
The residue is precisely the difference between the measurement and the story. The
process --- the linear story --- almost explains the measurement, and what is left when
it falls short is the residue, the second variation, the electron. The early demand,
left open until the paired machinery exists, is satisfied here by an identity.

The section's metaphysical anchor is that some quantities must be produced rather than
found, and that producing them is a measurement in its own right. The apparatus tanges
two coordinates at once into a single coupled measurement and funges the linear story
out of it by subtraction, and what remains is a signed unit that no single coordinate
could have shown. The electron is not discovered sitting in the vacuum; it is produced
by exciting the vacuum in a pair and reading what the linear story cannot cover. The
single invariant, asked to show itself on the simplest paired excitation, answers with
a charge.

\section{Strain if and only if Residue}

The book has used the word strain since its earliest chapters --- the informational
strain a variation leaves behind, the thing a process carries that its linear story
does not capture. This section identifies that strain with the residue, exactly, and
shows that the identification is an equivalence rather than a resemblance. The linear
story fails precisely when the residue is nonzero, and the strain is real precisely
when the charge is.

State the equivalence directly. The measurement and the story disagree if and only if
the residue is nonzero:
\[
  D(v) \neq r_{\mathrm{L}}(v_{\mathrm{L}}) + r_{\mathrm{R}}(v_{\mathrm{R}})
  \quad\Longleftrightarrow\quad
  \delta^2(v) \neq 0.
\]
This is immediate from the residue identity of the previous section --- the residue is
the difference between the two sides --- but its content is conceptual, not algebraic.
It says that the failure of the linear story and the presence of the residue are not
two phenomena that happen to coincide. They are one phenomenon. Wherever the story is
inadequate, a residue is present and measures exactly the inadequacy; wherever the
story is adequate, the residue is zero. Strain is the residue, with no gap between
them.

Make the strain concrete at the electron. The strain has two parts: a direction and a
failed response. The direction is that the pair moved --- both middle nodes left the
vacuum. The failed response is that the linear story did not account for the
measurement. At the electron both parts hold: the pair moved, and the story (four) did
not cover the measurement (three). So the strain is real:
\[
  \text{strain} = (\text{pair moved}) \wedge \neg(\text{story accounts for measurement}),
\]
and at the pair kick both conjuncts are true. The informational strain the book named
at the outset is not a metaphor here; it is a specific, checkable conjunction, and it
holds at the electron.

And it holds \emph{because} the electron is nonzero, and only because. The failure of
the story is equivalent to the charge being nonzero:
\[
  \neg(\text{story accounts for measurement})
  \quad\Longleftrightarrow\quad
  \text{electron} \neq 0.
\]
The strain is not an independent fact that happens to accompany the charge. It is the
charge, stated as a failure of the linear account. A zero charge would be a story that
holds, no strain; a nonzero charge is a story that fails, strain present. The electron
is real exactly to the degree that its residue refuses to vanish.

This lets the apparatus seat the electron into its own foundations without altering
them. The earliest chapters left a residue pressure inside the derivative's own
construction, but they did not yet have the paired machinery that could identify a
precise occupant. The electron's strain occupies that slot now, by name, and the
apparatus's own machinery manufactures the occupant: transmuting the paired variation
through the construction the early chapters already carried yields exactly the
electron's term. The earlier structure is unchanged; what was an unresolved pressure is
now placed in the slot the pair-production identity exposes.

The placement matters because it forecloses a suspicion. One might worry that the
electron was bolted onto the apparatus after the fact, an interpretation imposed on a
construction that had no pressure for it. The opposite is true. The construction left a
residue pressure that only the paired machinery can sharpen into a slot. The apparatus
does not bolt on a place for the electron; it finds, once the pair is moved, that its
own machinery produces both the place and the occupant together. The electron is native
to the structure, not a guest in it.

The section's metaphysical anchor is that strain and residue are the same thing read
two ways. The strain is the felt inadequacy of a linear story; the residue is the
measured remainder that inadequacy leaves; and they are equal, identically, at every
variation. The apparatus tanges the pair into a measurement, funges the story out of
it, and reads the strain as the residue that survives --- nonzero exactly when the
story fails, native to the slot exposed by the paired machinery. The book named the
strain before it could produce it. Here the strain is produced, named, and seated in
the place Ch12 can now identify.

\section{The C-Squared Certificate and Hilbert's Sixth, Finite Edition}

The configuration the pair kick excites can be certified, and the certificate has a
familiar name. A cubic spline carries data through its second derivative and nothing
above: it is twice continuously differentiable, what this section abbreviates
\(C^2\). The order-two data of such a spline is exactly the residue field the previous
sections produced, and the electron is its order-two quantum. This section states the
certificate, reads it against a century-old problem, and is careful about how much it
proves.

The certificate is a conjunction of three orders, supplied by a pair of certificates,
one per middle node. At order zero, both sub-paths admit --- the configuration is
constructible on each side of the middle. At order one, both Euler-Lagrange residuals
vanish in every direction --- the equations of motion hold. At order two, what remains
of every coupled variation is exactly the second variation --- the residue is the whole
of the order-two content, and its exactness leaves no order three to certify:
\[
  \underbrace{\text{admits left} \wedge \text{admits right}}_{\text{order } 0}
  \ \wedge\
  \underbrace{r_{\mathrm{L}} = 0 \wedge r_{\mathrm{R}} = 0}_{\text{order } 1}
  \ \wedge\
  \underbrace{D(v) = \delta^2(v)}_{\text{order } 2, \text{ exact}}.
\]
A configuration meeting all three is \(C^2\)-certified. Order two is the top, because
the order-two equality is exact: there is no remainder to push into a third order.

The word \(C^2\) deserves one honest gloss, because the finite apparatus cannot mean
by it quite what classical analysis means. Classical \(C^2\) asks the first and second
derivatives to be \emph{continuous}, a statement about a limit the finite corpus
cannot make. So the apparatus says what it can, and what it can say is in two respects
sharper. At order one it does not merely ask the first derivative to exist and be
continuous; it pins the first-order residual to zero, on-shell --- the equation of
motion holds, not just the derivative. And at order two it does not approximate; it
states an exact equality. \(C^2\) here means all the data through order two, with the
equation of motion holding and the second-order content exact. It is the only sense of
\(C^2\) the finite world can carry, and within that sense it is stronger than the
classical condition, not weaker.

The order-two data of the certified spline is the residue field, and so the electron
is the order-two quantum of the certificate. This is the unification the chapter has
been building toward. The certificate that says ``this configuration is smooth through
second order'' and the residue that says ``this is the electron's charge'' are the same
object read at the same order. The smoothness certificate carries, in its order-two
slot, exactly the signed unit the pair kick produces. On the vacuum --- the zero-cost
action --- that slot reads zero everywhere; the vacuum is empty, its residue identically
zero, and the certificate is inhabited there by the flat configuration the early
chapters established. The certificate is not conditional fiction; it holds on the
vacuum, with the residue field reading zero until a pair excites it.

Read against the history of mathematics, this is a small finite answer to a large old
question. In 1900 Hilbert asked, as his sixth problem, for the mathematical treatment
of the axioms of physics. What the apparatus offers is a variational mechanics with
equations of motion at order one, interaction at order two, and exactness at order
two, derived as a theorem from the apparatus's measurement primitives rather than
posited, and machine-checked. And the axioms it rests on are not the axioms of a
physical theory but two axioms of ordinary mathematics --- propositional extensionality
and the soundness of quotients --- with the physics derived from them and the
measurement primitives, not assumed. Two logical axioms, and a mechanics that follows.

The section is exact about how much this is, because the temptation to claim more is
real. It is one electron's worth of physics. There is no exclusion theorem here tying
the certificate to a vanishing electron, and there could not be: the certificate pins
first-order data, and the chapter's own central result is that first-order instruments
cannot see the coupling. Under a certificate the residue is not forced to zero --- it is
all that remains, by the order-two clause. What is established is at the level of
instances: the vacuum's residue reads zero everywhere, and the electron's action
refuses certification by an order-one obstruction --- its first variation does not
vanish --- not as a consequence of the electron. Hilbert asked for the axioms of all of
physics. This is one electron, certified through second order, on two logical axioms.

The section's metaphysical anchor is that smoothness and charge are the same order-two
fact. The apparatus tanges a configuration into a certificate that carries its data
through second order, and the order-two slot of that certificate is the residue field
the pair kick excites. To certify a configuration as \(C^2\) is to name its order-two
quantum, and that quantum, on a paired excitation of the vacuum, is the electron. The
oldest call for a mathematics of physics is answered here in the only currency the
apparatus trades in: a finite, exact, machine-checked certificate, one electron's
worth, resting on two axioms of ordinary logic.

\section*{Bridge: Where the Two Channels Come From}

The chapter produced a signed unit from a paired variation and certified the
configuration that carries it. In doing so it took one thing for granted: that there
were two channels to pair. The left side of the middle and the right side of the
middle were given as two coordinates, and the residue lived in the coupling between
them. But the apparatus never asked where the two channels came from. It moved a pair
because a pair was there to move.

That is the question the next chapter opens. The two-channel structure is not a
primitive the apparatus should help itself to. It is something the geometry must
generate. The energy of the earlier chapters and the boundary data of the loaded
anchor together carry more structure than the apparatus has yet drawn out, and that
structure includes the very split into two channels that pair production exploited.
The coupling has a geometry, and the geometry has a reason for there being two sides
to couple.

There is a related debt the next chapter must pay. The same single action has been
asked to do two jobs at once: to vanish on the vacuum, giving the empty residue, and
to produce a nonzero residue on an excitation. A vacuum channel and an excitation
channel are different demands, and the next chapter examines whether one action can
honestly serve both, or whether the geometry forces the two channels apart into
distinct roles that a single action cannot fill.

So the next chapter turns from the residue to its geometry. It shows how boundary
geometry --- the energy together with the depth of the boundary data --- generates the
two-channel gauge form, why the boundary's depth is a generative input rather than a
fixed frame, and why one action cannot do both the vacuum's job and the excitation's
job at once. The question the apparatus carries forward is the one its own success
hid: it produced the electron from two channels, but what produced the two channels?

\section*{Coda: The Depth That Needs Two Eyes}

Binocular vision performs this chapter's central fact with two eyes and a scene, and
it performs it so exactly that the mathematics can almost be read off the experience.
Depth, the quantity the scene does not flatly contain, is produced only by a pair, is
the difference between two flat stories, and carries a sign.

Begin with one eye. A single eye returns a flat image. It has color, shape, texture,
overlap --- a great deal of information --- but it does not, by itself, return depth.
Whatever cues a single eye exploits to guess at distance are inferences laid over a
fundamentally flat report; the eye's own measurement is two-dimensional. Close one eye
and reach for a thread through the eye of a needle, and the difficulty is immediate
and structural: the flat image does not carry the distance, and no amount of squinting
or attention recovers it. A one-eyed instrument is not poor at depth; it is, for the
geometric component of depth, blind by construction --- exactly as a one-coordinate
variation is blind to the coupling.

Open the second eye and depth appears, and it appears as a difference. The two eyes
see almost the same scene from slightly separated positions, and the small disparity
between their images --- the amount by which a near object shifts against the
background between the left view and the right view --- is what the visual system reads
as depth. Neither flat image contains the depth; the depth is what is left when the
two flat images are compared and the common part subtracted. Depth is the residue of
binocular vision: measurement by the pair, minus the story each eye tells alone.

And the residue is signed. An object nearer than the point the eyes are fixed on
shifts one way between the two images; an object farther than the fixation point shifts
the other way. The sign of the disparity distinguishes in front from behind, and the
two signs are not interchangeable --- they are the two orientations of one comparison,
exactly as the electron and the positron are the two orientations of one residue. The
visual system that reads a near object as near would read the same disparity with the
opposite sign as a far object: same magnitude, opposite sign, the pair the geometry
allows.

The lesson the eyes teach is the lesson of the pair kick. You cannot recover the depth
by improving one eye. A sharper monocular image is still flat; a brighter one is still
flat; the depth is not hiding in the single image waiting for better resolution,
because it was never a property of one image at all. It is a property of the pair, read
off the difference between two simultaneous flat reports. The apparatus that wanted the
residue had to move two coordinates at once for the same reason the visual system needs
two eyes open at once: the quantity lives in the comparison, not in either term.

That is the chapter in the body's own optics. One eye is the one-coordinate
instrument, flat and blind to the coupling; two eyes are the pair kick, producing what
neither could alone; the disparity is the residue, the difference between the
measurement and the flat stories; and the sign of the disparity is the difference
between the electron and the positron. The single invariant of the earlier chapters,
asked to show itself, shows itself the way depth does: only to an instrument willing to
open both eyes at once.
